\answer{ex:16:1}{血液是 $\tau=t_{1/2}/\ln2=14.4$~min的RC电路；最大值与最小值之比 $2^{T/t_{1/2}}=64$；基波衰减到 $0.55$；$t_{1/2}=T=60$~min时比值为 $2$。}
\answer{ex:16:2}{$\varphi^*=\arcsin\bigl((\omega-\omega_0)/K_\varphi\bigr)\approx41$~min（$\omega-\omega_0\approx0.047$~rad/d），弛豫时间 $1/(K_\varphi\cos\varphi^*)\approx3.9$~天；跳变 $+6$ 和 $-6$~h后分别需 $7.7$ 和 $7.9$~天。}
\answer{ex:16:3}{$K_c(n)=\sec^n(\pi/n)$：$8$ 和 $2.37$，误差 $11\%$ 和 $30\%$；$n\to\infty$ 时 $K_c\to1$，误差 $\to50\%$。}
\answer{ex:16:4}{极限指令（油门 $80$，刹车 $0$）经 $0.76$~s达到 $f=120$；“只用油门”需 $2.33$~s，慢三倍。}
\answer{ex:16:5}{$3\arctan(\omega\tau)+\omega d=\pi$，$K_c=(1+\omega^2\tau^2)^{3/2}$：$K_c=8$、$3.54$、$2.50$、$1.76$，周期 $3.63$、$5.46$、$6.86$、$9.27\,\tau$；$0.9K_c$ 时振荡衰减，$1.1K_c$ 时振荡稳定持续。}
\answer{ex:16:6}{$0.60\bar c$、$0.87\bar c$、$0.90\bar c$，极限 $0.91\bar c$；浓度恒定时响应为零：这样的接收方只读取剂量。}
\answer{ex:16:7}{开放问题。在反馈机制下，只有 $K>8$ 时才有节律，周期正比于各级的 $\tau$，几乎与 $K$ 无关（$K=8.5$ 时 $3.64\,\tau$，$K=40$ 时 $4.23\,\tau$）：$K$ 改变一点五倍只使周期偏移 $2$--$4\%$，小于测量误差。能作出判决的是断开环路（在第一级输入端维持恒定的终端激素会消灭这种节律），或者观察在周期不同相位加入少量激素后的响应。}
